% Copyright 2026 MarcosHCK
%
\documentclass{article}[12pt,a4paper]
\usepackage{amsfonts}
\usepackage{amsmath}
\usepackage{amsthm}
\usepackage{bookmark}
\usepackage{hyperref}
\usepackage{cleveref} % must load after hyperref
\usepackage{graphicx}
\usepackage{inputenc}[utf8]
\usepackage{indentfirst}
\usepackage{polyglossia}[spanish,english]
\usepackage{biblatex}[style=apa] % must load after polyglossia
\usepackage{csquotes} % must load after biblatex
\usepackage{setspace}
\usepackage{svg}
\usepackage{xcolor}[dvipsnames]

\setmainlanguage{spanish}
\setotherlanguage{english}

\setstretch{1.5}

\newcommand{\myeoch}{\pagebreak}

\definecolor{Aquamarine}{rgb}{0.498, 1, 0.83}
\definecolor{DarkAquamarine}{rgb}{0.18, 0.55, 0.34}
\definecolor{YellowOrange}{rgb}{1, 0.7, 0}
\newcommand\myshade{85}
\colorlet{mycitecolor}{YellowOrange}
\colorlet{mylinkcolor}{violet}
\colorlet{myurlcolor}{DarkAquamarine}

\hypersetup {
    colorlinks = true,
    citecolor  = mycitecolor!\myshade!black,
    linkcolor  = mylinkcolor!\myshade!black,
    urlcolor   = myurlcolor!\myshade!black,
  }

\begin{document}

\begin{titlepage}

  \centering
  {\huge\bfseries Reporte \#1\par}
	{\Large Autor: \itshape Marcos Antonio Pérez Lorenzo\par}
\end{titlepage}

\pagebreak

\section{Teoría}

\subsection{High-pass}
Sea el circuito RC de la forma:

\begin{proof}

  El voltaje de entrada a través del tiempo se expresa como (usando la ley de Kirchhof):

  $$V_{i}(t)=V_{r}(t)+V_{c}(t)$$

  Si se expanden las formulas correspondientes a las voltajes:

  $$V_{i}(t)=I(t)*R+\frac{1}{C}\int I(t)\,dt$$

  Utilizando la transformada de Laplace:

  $$V_{i}(s)=I(s)*R+\frac{1}{C}\frac{1}{s}*I(s)$$

  \begin{equation}
    V_{i}(s)=I(s)\left(R+\frac{1}{Cs}\right)
    \label{eq:high1}
  \end{equation}

  Para el voltaje de salida es similar:

  $$V_{o}(t)=V_{r}(t)$$
  $$V_{o}(t)=I(t)*R$$

  \begin{equation}
    V_{o}(s)=I(s)*R
    \label{eq:high2}
  \end{equation}

  Finalmente, para calcular la expresión de la transformación \(T(s)\), se combinan \ref{eq:high1} y \ref{eq:high2}:

  $$T(s)=\frac{V_{o}(s)}{V_{i}(s)}$$
  $$T(s)=\frac{I(s)R}{I(s)\left(R+\frac{1}{Cs}\right)}$$
  $$T(s)=\frac{R}{R+\frac{1}{Cs}}$$
  $$T(s)=\frac{1}{1+\frac{1}{RCs}}$$
  $$T(s)=\frac{1}{\frac{1+RCs}{RCs}}$$

  \begin{equation}
    T(s)=\frac{RCs}{1+RCs}
    \label{eq:rhigh}
  \end{equation}
\end{proof}

\subsection{Low-pass}
Sea el circuito RC de la forma:

\begin{proof}

  El voltaje de entrada a través del tiempo se expresa como (usando la ley de Kirchhof):

  $$V_{i}(t)=V_{r}(t)+V_{c}(t)$$

  Si se expanden las formulas correspondientes a las voltajes:

  $$V_{i}(t)=I(t)*R+\frac{1}{C}\int I(t)\,dt$$

  Utilizando la transformada de Laplace:

  $$V_{i}(s)=I(s)*R+\frac{1}{C}\frac{1}{s}*I(s)$$

  \begin{equation}
    V_{i}(s)=I(s)\left(R+\frac{1}{Cs}\right)
    \label{eq:low1}
  \end{equation}

  Para el voltaje de salida es similar:

  $$V_{o}(t)=V_{c}(t)$$
  $$V_{o}(t)=\frac{1}{C}\int I(t)\,dt$$

  Aplicando la transformación de Laplace nuevamente:

  $$V_{o}(s)=\frac{1}{C}\frac{1}{s}I(s)$$

  \begin{equation}
    V_{o}(s)=\frac{I(s)}{Cs}
    \label{eq:low2}
  \end{equation}

  Finalmente, para calcular la expresión de la transformación \(T(s)\), se combinan \ref{eq:low1} y \ref{eq:low2}:

  $$T(s)=\frac{V_{o}(s)}{V_{i}(s)}$$
  $$T(s)=\frac{\frac{I(s)}{Cs}}{I(s)\left(R+\frac{1}{Cs}\right)}$$
  $$T(s)=\frac{1}{Cs*\left(R+\frac{1}{Cs}\right)}$$

  \begin{equation}
    T(s)=\frac{1}{1+RCs}
    \label{eq:rlow}
  \end{equation}
\end{proof}

\section{Casos}

\subsection{Dominio de frecuencias}

Para representar ambas ecuaciones \ref{eq:rhigh} y \ref{eq:rlow}, debemos primero transformar ambas de funciones complejas a magnitudes reales, substituyendo (donde \(f\) es la frecuencia):

\begin{equation*}
  s=\jmath\omega=\jmath2\pi f
\end{equation*}

Luego, para el filtro high-pass:

\begin{proof}

  $$T(\jmath\omega)=\frac{RC\jmath\omega}{1+RC\jmath\omega}$$
  $$\left|T(\jmath\omega)\right|=\frac{\left|RC\jmath\omega\right|}{\left|1+RC\jmath\omega\right|}$$

  \begin{equation}
    \left|T(\jmath\omega)\right|=\frac{\omega RC}{\sqrt{1+\left(\omega RC\right)^2}}
    \label{eq:realhigh}
  \end{equation}
\end{proof}

En el caso del filtro low-pass:

\begin{proof}

  $$T(\jmath\omega)=\frac{1}{1+RC\jmath\omega}$$
  $$\left|T(\jmath\omega)\right|=\frac{1}{\left|1+RC\jmath\omega\right|}$$

  \begin{equation}
    \left|T(\jmath\omega)\right|=\frac{1}{\sqrt{1+\left(\omega RC\right)^2}}
    \label{eq:reallow}
  \end{equation}
\end{proof}

\subsection{Gráfica}

\begin{figure*}
  \includesvg[width=1\textwidth]{high_pass.svg}
  \caption{Gráfica para el filtro high-pass}
\end{figure*}

\begin{figure*}
  \includesvg[width=1\textwidth]{low_pass.svg}
  \caption{Gráfica para el filtro low-pass}
\end{figure*}

\end{document}